% Public text-only snapshot of the instructor's Module 12 diffusion article.
% Upstream: /Users/setup/Library/CloudStorage/Box-Box/Teaching/6050/LaTeX/Module 12/12.diffusion.tex
% Upstream SHA-256: 261a009143acc821da578a4480b48d45cb584df2779a19772c20d97bd714bcf4
% TikZ source, executable code, external-library workflows, runtime downloads,
% product-current surveys, hardware claims, external assets, and assignment-specific
% exercises are omitted. The retained lecture spine and equations are provenance
% records, not endorsed empirical results; the book independently checks and derives
% every published claim, implementation, experiment, and figure.
\documentclass{article}
\usepackage{amsmath}
\usepackage{amssymb}
\usepackage[margin=1in]{geometry}
\usepackage{booktabs}
\usepackage{hyperref}

\title{Diffusion Models and Conditional Generation}
\author{Deep Learning Class}
\date{}

\begin{document}
\maketitle

\section*{Public Snapshot Boundary}

This text-only snapshot preserves the instructor-authored conceptual sequence and
selected equations. Diagram source, code, model-loading examples, product comparisons,
hardware prescriptions, and claims tied to a particular software or product release
have been removed. Technical claims must be verified against primary sources before
book publication.

\section{The Generative Modeling Landscape}

The generative objective is to learn a model distribution $p_{\theta}(\mathbf{x})$
that approximates a data distribution. Depending on the model family, the learned
object may support sampling, likelihood or likelihood-bound evaluation, and
conditional generation $p_{\theta}(\mathbf{x}\mid\mathbf{c})$.

An autoregressive model factors a joint distribution as

\[
p_{\theta}(\mathbf{x})
=\prod_{t=1}^{T}p_{\theta}(x_t\mid\mathbf{x}_{<t}).
\]

A generative adversarial model trains a generator against a discriminator. A diffusion
model instead defines a fixed corruption process and learns a reverse denoising
process.

% [The original paradigm table, named-product examples, state-of-the-art claims, and
% all TikZ diagrams are omitted.]

\section{DDPM: Forward Noising}

For data $\mathbf{x}_0\sim p_{\mathrm{data}}$, the forward process adds Gaussian
noise according to

\[
q(\mathbf{x}_t\mid\mathbf{x}_{t-1})
=\mathcal{N}\!\left(
\mathbf{x}_t;
\sqrt{1-\beta_t}\,\mathbf{x}_{t-1},
\beta_t I
\right).
\]

Equivalently,

\[
\mathbf{x}_t
=\sqrt{1-\beta_t}\,\mathbf{x}_{t-1}
+\sqrt{\beta_t}\,\boldsymbol{\epsilon},
\qquad
\boldsymbol{\epsilon}\sim\mathcal{N}(0,I).
\]

Define $\alpha_t=1-\beta_t$ and
$\bar\alpha_t=\prod_{s=1}^{t}\alpha_s$. Gaussian composition gives the direct
marginal

\[
q(\mathbf{x}_t\mid\mathbf{x}_0)
=\mathcal{N}\!\left(
\mathbf{x}_t;
\sqrt{\bar\alpha_t}\,\mathbf{x}_0,
(1-\bar\alpha_t)I
\right),
\]

or

\[
\mathbf{x}_t
=\sqrt{\bar\alpha_t}\,\mathbf{x}_0
+\sqrt{1-\bar\alpha_t}\,\boldsymbol{\epsilon}.
\]

This closed form permits training at a randomly selected timestep without simulating
all preceding corruption steps.

\section{The Learned Reverse Process}

Generation begins with a noise sample and repeatedly models
$p_{\theta}(\mathbf{x}_{t-1}\mid\mathbf{x}_t)$. A Gaussian parameterization is

\[
p_{\theta}(\mathbf{x}_{t-1}\mid\mathbf{x}_t)
=\mathcal{N}\!\left(
\mathbf{x}_{t-1};
\boldsymbol{\mu}_{\theta}(\mathbf{x}_t,t),
\boldsymbol{\Sigma}_{\theta}(\mathbf{x}_t,t)
\right).
\]

One common parameterization fixes the reverse variance and predicts the noise used in
the forward marginal. The simplified training objective is

\[
\mathcal{L}_{\mathrm{simple}}
=\mathbb{E}_{t,\mathbf{x}_0,\boldsymbol{\epsilon}}
\left[
\left\|
\boldsymbol{\epsilon}
-\boldsymbol{\epsilon}_{\theta}(\mathbf{x}_t,t)
\right\|_2^2
\right],
\]

where

\[
\mathbf{x}_t
=\sqrt{\bar\alpha_t}\,\mathbf{x}_0
+\sqrt{1-\bar\alpha_t}\,\boldsymbol{\epsilon}.
\]

The corresponding mean parameterization in the lecture is

\[
\boldsymbol{\mu}_{\theta}(\mathbf{x}_t,t)
=\frac{1}{\sqrt{\alpha_t}}
\left(
\mathbf{x}_t
-\frac{\beta_t}{\sqrt{1-\bar\alpha_t}}
\boldsymbol{\epsilon}_{\theta}(\mathbf{x}_t,t)
\right).
\]

Sampling starts from $\mathbf{x}_T\sim\mathcal{N}(0,I)$ and applies the learned
reverse transition for $t=T,T-1,\ldots,1$, retaining the prescribed stochastic term
where appropriate.

\section{Architecture and Time Conditioning}

The denoiser receives both a noisy sample $\mathbf{x}_t$ and its timestep $t$.
The lecture uses a multi-resolution encoder--decoder with skip connections as the
architectural example and injects a timestep embedding throughout the network. A
sinusoidal representation has the schematic form

\[
t_{\mathrm{emb}}
=\left[
\sin(\omega_1t),\cos(\omega_1t),
\sin(\omega_2t),\cos(\omega_2t),\ldots
\right].
\]

% [The U-Net diagram is omitted.]
% [The complete implementation section is omitted. Its original simplified U-Net
% listing also contains incompatible skip-connection shapes and must not be ported.]

\section{Latent Diffusion}

Latent diffusion separates perceptual compression from generative modeling. An encoder
$\mathcal{E}$ maps observations to a latent representation, diffusion is trained in
that representation, and a decoder $\mathcal{D}$ maps generated latents back to the
observation space:

\[
\mathbf{z}_0=\mathcal{E}(\mathbf{x}),
\qquad
\hat{\mathbf{x}}=\mathcal{D}(\mathbf{z}_0).
\]

This changes the state space on which the denoiser operates. It does not by itself
establish a particular wall-clock speedup, memory requirement, or fidelity guarantee.

% [Named-product pipelines, release matrices, hardware claims, and model-loading code
% are omitted.]

\section{Transformer and Text Conditioning}

A diffusion transformer tokenizes a noisy latent, adds position and timestep
information, applies transformer blocks, and maps the result back to the latent grid.
For text conditioning, the lecture uses cross-attention:

\[
\operatorname{Attention}(Q,K,V)
=\operatorname{softmax}\!\left(\frac{QK^{\top}}{\sqrt d}\right)V,
\]

with queries from the noisy image representation and keys and values from the text
representation.

Classifier-free guidance trains conditional and unconditional predictions. Under the
lecture's convention, the guided prediction is

\[
\tilde{\boldsymbol{\epsilon}}
=(1+s)\boldsymbol{\epsilon}_{\theta}(\mathbf{x}_t,t,\mathbf{c})
-s\boldsymbol{\epsilon}_{\theta}(\mathbf{x}_t,t,\varnothing).
\]

The scale convention must be stated because an alternative notation writes the same
interpolation or extrapolation with a differently shifted coefficient.

% [External text-encoder code, library APIs, and product-specific conditioning claims
% are omitted.]

\section{Evaluation and Responsibility}

The lecture distinguishes sample-quality, distributional, conditioning, and human
evaluation. It also raises dataset bias, consent, provenance, misuse, and watermarking
as evaluation and governance concerns.

% [Metric-library code and claims that any one score certifies quality are omitted.]

\section{Primary Reading Named in the Lecture}

\begin{itemize}
\item Ho et al., ``Denoising Diffusion Probabilistic Models'' (2020).
\item Song et al., ``Score-Based Generative Modeling through Stochastic Differential
Equations'' (2021).
\item Song et al., ``Denoising Diffusion Implicit Models'' (2021).
\item Rombach et al., ``High-Resolution Image Synthesis with Latent Diffusion Models''
(2022).
\item Peebles and Xie, ``Scalable Diffusion Models with Transformers'' (2023).
\item Ho and Salimans, ``Classifier-Free Diffusion Guidance'' (2022).
\end{itemize}

% [Product-current resource lists, lab time estimates, and assignment exercises are
% omitted.]

\end{document}
